\section{Numerical experiments}
\label{sec:experiments}

We compare \texttt{.tmu} with \texttt{.tex} on compiling/rendering speed, LLM task performance, and fine-tuning.

\paragraph{Setup.} All experiments derive from a 120-page arXiv document (\texttt{arXiv:2502.17655}) and run on an Intel Ultra 9 285H CPU, GeForce RTX 5080 GPU, and 32GB RAM. We evaluate four LLMs---Deepseek-v3.2, Claude-Sonnet-4.5, Gemini-3-pro, and ChatGPT-5.2---on structure locating (20 questions), style-aware file merging (5 tasks), and debugging (20 samples). Scores reward correctness while penalizing total token usage $T$ (input, thinking, output, and MCP tools), and timing benchmarks report the average of three trials per document. We refer to these scores as Cost-aware Task Utility. Incorrect outputs receive zero, while token cost discounts successful completions; a correct but over-budget response may also score zero, without being considered semantically equivalent to a hallucination.

\subsection{Benchmark on compiling/rendering time}

We benchmark compilation and rendering on 6 arXiv documents. \texttt{.tmu} outperforms \texttt{.tex} on both full compilation and incremental updates; the latter is most decisive (Figure~\ref{fig:inc-update}), because \texttt{.tmu} updates local tree nodes without I/O overhead ($t_{\text{IO}}=0$, activated only when opening the document). These low feedback delays directly reduce agent time cost in real-time LLM writing. \emph{(Full-compilation results and the $t_{\text{IO}}$ cost model are in Appendix~\ref{app:bench}.)}

\begin{figure}[htbp]
    \centering
    \includegraphics[width=\columnwidth]{figure/inc-update.pdf}
    \caption{Benchmark on incremental update. Each bar represents the average of three trials.}
    \label{fig:inc-update}
\end{figure}

\subsection{Performance in LLM tasks}

The standardized grammar and explicit tree-node tags in \texttt{.tmu} help models locate targets, merge contexts, and debug ill-formed structures more efficiently than \texttt{.tex}. Figure~\ref{fig:llm-tasks} summarizes the three LLM task results.

\begin{figure*}[t]
    \centering
    \begin{subfigure}[t]{0.32\textwidth}
        \centering
        \includegraphics[width=\linewidth]{figure/reading.pdf}
        \caption{Structure locating}
        \label{fig:reading}
    \end{subfigure}\hfill
    \begin{subfigure}[t]{0.32\textwidth}
        \centering
        \includegraphics[width=\linewidth]{figure/writing.pdf}
        \caption{Style-aware merging}
        \label{fig:writing}
    \end{subfigure}\hfill
    \begin{subfigure}[t]{0.32\textwidth}
        \centering
        \includegraphics[width=\linewidth]{figure/debugging.pdf}
        \caption{Debugging}
        \label{fig:debugging}
    \end{subfigure}
    \caption{LLM task performance for \texttt{.tmu} and \texttt{.tex}.}
    \label{fig:llm-tasks}
\end{figure*}

\subsubsection{Locating document structure}

To evaluate structural understanding, each LLM answers 20 document-structure questions (Appendix~\ref{app:20ques}). For each answer, we take
\begingroup\small
\[
    u_s = \max\left(0,\ \begin{cases}
        5 - \left\lfloor \dfrac{T}{1 \times 10^4} \right\rfloor & \text{right answer,} \\[6pt]
        0 & \text{wrong answer,}
    \end{cases} \right)
\]
\endgroup
where $\sum u_s \in [0, 100]$. The 10k-token scale separates confident structure location from longer reasoning/tool-use loops, which indicate lower efficiency and higher cost~\cite{ling_table2latex_2025}.

Despite LLMs' stronger pretraining exposure to \LaTeX{}, \texttt{.tmu} leads for most models (Figure~\ref{fig:reading}). Its reference metadata is updated and stored directly in the file (e.g., \path!<associate|sec:tree-struc-on-mogan|<tuple|5.1|13>>!), whereas \texttt{.tex} requires scanning and resolving the compiled numbering context.

Gemini is the only exception, scoring 90 on \texttt{.tex} versus 87 on \texttt{.tmu}. Thus, we claim gains in most, rather than all, settings and interpret this as a comparison of the complete native representations rather than tree syntax alone.
% Replaced by the combined Figure~\ref{fig:llm-tasks} to save space.
% \begin{figure}[htbp]
%     \centering
%     \includegraphics[width=\columnwidth]{figure/reading.pdf}
%     \caption{Test on locating document structure.}
%     \label{fig:reading}
% \end{figure}

\subsubsection{Merging files with distinct doc-styles}
\label{sec:doc-style}

We define \emph{doc-style} as macro naming and command usage. Documents may use different theorem aliases, redefine the same macro differently, or give the same command incompatible meanings.

We ask LLMs to complete two assignments.

Assignment 1 generates two independently compilable files: \texttt{theorems.tex}, a theorem collection with many new macros and packages, and \texttt{proofs.tex}, a disordered proof collection with a distinct doc-style and conflicting packages. The files are connected by theorem and equation cross-references.

Assignment 2 merges the generated files into the doc-style of \texttt{theorems.tex}, while preserving compilability and placing each proof below the referenced theorem.

We obtain \texttt{theorems.tmu} and \texttt{proofs.tmu} from their \texttt{.tex} versions; rendered content is identical, as guaranteed by a \LaTeX{} importing engine.

Assignment 1 is 1 generation task; Assignment 2 is 4 merge tasks. For each task, we take
\begin{equation*}
\resizebox{\columnwidth}{!}{$
u_m =
\max\left(0,\
\begin{cases}
20 - 2E_{\mathrm{ref}} - \left\lfloor \dfrac{T}{10^4} \right\rfloor - E_{\mathrm{sty}}
    & \text{success on the first try,} \\[6pt]
10 - 2E_{\mathrm{ref}} - \left\lfloor \dfrac{T}{10^4} \right\rfloor - E_{\mathrm{sty}}
    & \text{success on the second try,} \\[6pt]
0
    & \text{fail within two tries,}
\end{cases}
\right)
$}
\end{equation*}
where $E_{\text{ref}}$ is the number of reference failures (i.e., "??" appears but compilation succeeds), $E_{\text{sty}}$ is the number of cases where the merged file has the doc-style of \texttt{proofs.tex} (we required LLMs to write the merged file in the doc-style of \texttt{theorems.tex}, so other macro aliases or redefined macros are not accepted), $\sum u_m \in [0, 100]$.

As shown in Figure~\ref{fig:writing}, \texttt{.tmu} scores higher for all LLMs. Its grammatical consistency removes many macro-conflict and style-unification choices that make complex \LaTeX{} generation brittle~\cite{kale2025texpert}; merging structured contexts becomes closer to copying compatible subtrees.
% Replaced by the combined Figure~\ref{fig:llm-tasks} to save space.
% \begin{figure}[htbp]
%     \centering
%     \includegraphics[width=\columnwidth]{figure/writing.pdf}
%     \caption{Test on merging files with distinct doc-styles.}
%     \label{fig:writing}
% \end{figure}
%
% Furthermore, with the \texttt{.tmu} representation, merging contexts from two documents with distinct doc-styles reduces to a simple copy-and-paste task.

\subsubsection{Debugging ill-formed documents using error messages}

Debugging ill-formed documents is a common LLM co-writing use case. We construct 20 samples from the source document (Table~\ref{tab:ill-dist}), feed error messages to LLMs, and ask them to repair the document.

\begin{table}[t]
    \centering
    \small
    \setlength{\tabcolsep}{4pt}
    \caption{Distribution of error types.}
    \label{tab:ill-dist}
    \begin{tabular}{lc}
        \toprule
        \textbf{Error Type} & \textbf{Count} \\
        \midrule
        Unclosed bracket & 4 \\
        Unclosed environment & 5 \\
        Wrong command usage & 4 \\
        Undefined cross-reference & 3 \\
        Conflicting packages & 2 \\
        Self-recursive macros & 2 \\
        \bottomrule
    \end{tabular}
\end{table}

% Older full table, kept for restoration if needed.
% \begin{table}[htbp]
%     \centering
%     \small
%     \setlength{\tabcolsep}{3pt}
%     \caption{Distribution of error types in samples}
%     \label{tab:ill-dist}
%     \begin{tabular}{lc}
%         \toprule
%         \textbf{Error Type}     & \textbf{Count} \\
%         \midrule
%         Unclosed bracket          & 4              \\
%         Unclosed environment      & 5              \\
%         Wrong command usage       & 4              \\
%         Undefined cross-reference & 3              \\
%         Conflicting packages      & 2              \\
%         Self-recursive macros     & 2              \\
%         \bottomrule
%     \end{tabular}
% \end{table}

For each error, we take
\begingroup\small
\[
    u_d = \max\left(0,\ \begin{cases}
        5 - \left\lfloor \dfrac{T}{1 \times 10^4} \right\rfloor & \text{right answer,} \\[6pt]
        0 & \text{wrong answer,}
    \end{cases} \right)
\]
\endgroup
where $\sum u_d \in [0, 100]$.

 We treat this as an end-to-end evaluation under each format's native validation and diagnostic workflow, rather than a controlled comparison using identical plain-text error messages.

The \texttt{.tex} group requires substantially more reasoning on most samples, while the \texttt{.tmu} group solves all samples and only two exceed 10k tokens (Figure~\ref{fig:debugging}). A defect in \texttt{.tmu} affects only the closest tree-tag ancestor (Section~\ref{sec:tree-struc-on-mogan}), yielding localized messages instead of long logs detached from root causes.
% Replaced by the combined Figure~\ref{fig:llm-tasks} to save space.
% \begin{figure}[htbp]
%     \centering
%     \includegraphics[width=\columnwidth]{figure/debugging.pdf}
%     \caption{Test on debugging ill-formed documents using error messages.}
%     \label{fig:debugging}
% \end{figure}

Well-formed \texttt{.tmu} also precludes classes such as unclosed environments or self-recursive macros; when a malformed tree exists, a structured editor can surface the affected node directly while the rest of the document remains renderable.

\subsection{Efficiency in fine-tuning}
\label{sec:eff-in-sft}

Because \texttt{.tmu} is tree-structured while \texttt{.tex} follows a linear macro flow, next-token prediction should be less ambiguous in \texttt{.tmu}.

We fine-tune Qwen2.5-7B-Instruct~\cite{qwen2025qwen25} with LoRA on 1000 random formulas, converting each \LaTeX{} formula to an equivalent rendered \texttt{.tmu} S-expression. The formulas cover common mathematical constructs; each is split into a prefix prompt and completion target.

Following Dong et al.~\cite{dong2025machinelearninglm}, both runs use identical LoRA settings and differ only in training format. Figure~\ref{fig:fine-tune} shows that \texttt{.tmu} converges to around 0.4 loss, while \texttt{.tex} converges to around 0.7. This supports the low-entropy claim: equivalent \LaTeX{} strings such as \verb|\frac{a}{b}| and \verb|{a \over b}|, or \verb|x^{2}| and \verb|x^2|, create extra prediction uncertainty absent from canonicalized \texttt{.tmu} trees.

\begin{figure}[htbp]
    \centering
    \includegraphics[width=\columnwidth]{figure/fine-tune.pdf}
    \caption{Experiment of LoRA based on Qwen2.5-7B-Instruct.}
    \label{fig:fine-tune}
\end{figure}

This advantage matters especially for large documents with multiple doc-styles, where \LaTeX{} models must infer commands from preamble-local macro definitions~\cite{lin2024accurate}.

Moreover, discussions of the \texttt{.tmu} format versus Markdown are attached in Limitations.
